% --- Page Style ---
\pagestyle{fancy}
\fancyhf{}
\lhead{Course: CS 2110 - Discrete Mathematics}
\rhead{Exam ID: \underline{\hspace{2cm}}}
\lfoot{Georgia Institute of Technology}
\rfoot{Page \thepage}
\renewcommand{\headrulewidth}{0.4pt}
\renewcommand{\footrulewidth}{0.4pt}

% --- Header spanning both columns ---
\twocolumn[
  \begin{center}
    {\huge \textbf{Midterm Examination II}} \\
    \vspace{0.2cm}
    {\large \textbf{Topics: Number Theory and Abstract Algebra}} \\
    \vspace{0.4cm}
    \begin{tabular}{|l|l|}
      \hline
      \textbf{Student Name:} & \hspace{6cm} \\ \hline
      \textbf{GTID:} & \hspace{6cm} \\ \hline
    \end{tabular}
  \end{center}
  \vspace{0.6cm}
  \textbf{Instructions:} Answer all questions in the space provided. You may use the back of the pages for scratch work. Calculators are not permitted.
  \vspace{0.8cm}
]

% --- Questions ---

\section*{Part I: Number Theory}

\begin{enumerate}[leftmargin=*, label=\textbf{\arabic*.}]

  \item Let $n \in \mathbb{Z}^+$. Define the reduced residue system modulo $n$ as $Z_n^x$. If $n = 12$, list all elements of $Z_n^x$. \points{5}
  \vspace{2cm}

  \item Prove Wilson's Theorem: For any prime $p$, $(p-1)! \equiv -1 \pmod{p}$. \points{10}
  \vspace{4cm}

  \item Solve the following system of linear congruences using the Chinese Remainder Theorem:
  \begin{align*}
    x &\equiv 2 \pmod{3} \\
    x &\equiv 3 \pmod{5} \\
    x &\equiv 2 \pmod{7}
  \end{align*} \points{10}
  \vspace{4cm}

\end{enumerate}

\section*{Part II: Abstract Algebra}

\begin{enumerate}[leftmargin=*, label=\textbf{\arabic*.}, resume]

  \item Define a group $(G, \ast)$. Prove that the identity element $e \in G$ is unique. \points{5}
  \vspace{3cm}

  \item Let $H$ be a subgroup of $G$. Define the left cosets of $H$ in $G$ and state Lagrange's Theorem regarding the order of $H$. \points{10}
  \vspace{4cm}

  \item \textbf{True/False:} Every cyclic group is Abelian. Briefly justify your answer. \points{5}
  \vspace{2cm}

\end{enumerate}

\section*{Part III: Bonus}

\begin{enumerate}[leftmargin=*, label=\textbf{\arabic*.}, resume]
  \item Find the greatest lower bound (infimum) and least upper bound (supremum) for the set $S = \{ \frac{1}{n} : n \in \mathbb{Z}^+ \}$. \points{5}
\end{enumerate}
